\scanpage{9}\mspage{6}
\begingroup
\makeatletter\tagsleft@true\makeatother
5. \uline{Application}: \uline{The} \uline{relation} \uline{between} $L_v(1)$ \uline{and} $\int_{E_v}|\omega_v|$ \uline{in} \uline{case} $k$ \uline{is} \uline{finite}.

Suppose now $R$ is complete and $k$ is finite with $q=\operatorname{Card} k$. On $K$ we agree to use the additive Haar measure with respect to which $R$ has measure $1$. This being agreed, a differential \inserted{$\omega$} on $E$ gives us a measure $|\omega|$ on $E(K)$ in the usual manner.

\uline{Corollary 5.1} (\unsure{of|cf} theorem 4.2) $\displaystyle\int_{E_1(K)}|\omega|=\frac{1}{q}$, if $\omega$ \uline{is} \uline{a} \uline{differential} \uline{of} $1^{\mathrm{st}}$ \uline{kind} \uline{on} $E$ \uline{coming} \uline{from} \uline{a} \uline{minimal} \uline{equation}.

Indeed, by (4.2) we have $\displaystyle\int_{E_1(K)}|\omega|=\int_{\mathfrak p}|dz|=\frac{\int_R|dz|}{(R:\mathfrak p)}=\frac{1}{q}$

The local factor occuring in the Euler product with a good functional equation is (Serre tells me) as follows, in which $N_v=\operatorname{Card}(\widetilde E_0(k))$\unsure{:|;|.}
{\small
\[
L_v(s)=\left\{
\begin{array}{@{}c@{\quad}l@{}}
\dfrac{1}{1-(q+1-N_v)q^{-s}+q^{1-2s}}
 & \parbox{.39\linewidth}{, if $\widetilde E$ non singular}\\[1.0ex]
\dfrac{1}{1-q^{-s}}
 & \parbox{.39\linewidth}{, if $\widetilde E$ has a node with the two tangents ratl over $k$ (in which case $N_v=q-1$, and $\widetilde E_0$ is the multiplicative group.)}\\[1.0ex]
\dfrac{1}{1+q^{-s}}
 & \parbox{.39\linewidth}{, if $\widetilde E$ has a node with irrational tangents (in which case $N_v=q+1$, and $\widetilde E_0$ is the twisted mult. group)}\\[1.0ex]
1 & \parbox{.39\linewidth}{, if $\widetilde E$ has a cusp (in which case $N_v=q$, and $\widetilde E_0$ is the additive gp.)}
\end{array}\right.
\tag*{(5.2)}
\]
}

In all cases therefore, $L_v(1)=\dfrac{q}{N_v}$. Since $N_v=(E_0(K):E_1(K))$: it follows from cor 5.1 that we have $\displaystyle\int_{E_0(K)}|\omega|=(L_v(1))^{-1}$, because $|\omega|$ is invariant under translation. Finally then,

\uline{Theorem 5.2}: \uline{If} \uline{we} \uline{use} \uline{the} \uline{measure} \uline{on} $K$ \uline{for} \uline{which} $R$ \uline{gets} \uline{measure} $1$, \uline{and} \uline{use} \uline{a} \uline{differential} \uline{of} \uline{first} \uline{kind} $\omega$ \uline{coming} \uline{from} \uline{a} \uline{minimal} \uline{equation}, \uline{then}
\[
\int_{E(K)}|\omega|=\frac{(E(K):E_0(K))}{L_v(1)}
\tag*{(5.3)}
\]
In other words, the ``fudge factors''\ of Birch and Sw-Dyer are just the indices $(E(K):E_0(K))$,
\endgroup

% archive sheet number: 8
